\begin{table}[htbp]
\centering
\caption{External Validity Scope: Modality and Commodity-Service Boundaries}
\label{tab:external_validity_scope}
\begin{threeparttable}
\small
\begin{tabular}{lllcc}
\toprule
Dimension & Group & Metric & Value & 95\% CI \\
\midrule
coverage & Commodity & share\_items & $0.887$ & --- \\
coverage & Service & share\_items & $0.113$ & --- \\
dominant\_item\_class & Commodity & auc\_log\_tc & $0.905$ & $[0.879, 0.931]$ \\
dominant\_item\_class & Service & auc\_log\_tc & $0.966$ & $[0.961, 0.971]$ \\
item\_class\_price & Commodity & coef\_losers & $0.045$ & $[0.002, 0.088]$ \\
item\_class\_price & Service & coef\_losers & $0.128$ & $[0.024, 0.231]$ \\
modal\_primary\_auc & convite\_primary & auc\_log\_tc & $0.816$ & $[0.758, 0.874]$ \\
modal\_primary\_auc & pregao\_primary & auc\_log\_tc & $0.952$ & $[0.946, 0.958]$ \\
\bottomrule
\end{tabular}
\begin{tablenotes}
\small
\item \textit{Notes:} Item classes use a transparent text rule on BEC item descriptions. Rows labelled \texttt{modal\_primary\_auc} reuse the always-loser firm split by primary modality; rows labelled \texttt{item\_class\_price} report the within-item price association separately for commodity-like and service-like items. The point is not to claim exportability to all procurement, but to document that the paper speaks most naturally to standardized-goods environments and selected simple services.
\item \textit{Source:} \texttt{scripts/52\_external\_validity\_scope.R}.
\end{tablenotes}
\end{threeparttable}
\end{table}
